\chapter{第一讲\quad 标准波动方程的估计}

考虑 Minkowski 时空 $\mathbb{R}^{n+1}$ 中的标准波动方程
\begin{equation}
\Box\phi=0.\end{equation}

Minkowski 空间 $\mathbb{R}^{n+1}$ 中的标准惯性坐标记为 $x^\mu$, $\mu = 0,1,\ldots,n$, 相对这组坐标 Minkowski 度量取对角形式 $m_{\mu\nu}= \mathrm{diag}(-1,1,\ldots,1)$. 我们有 $x^0 = t$, $x = (x^1,\ldots,x^n)$ 表示空间坐标. 下文中我们始终沿用关于重复指标求和的约定, 以及关于向量与张量指标升降的标准约定. 特别地, 若 $x_\mu = m_{\mu\nu}x^\nu$, 则 $x_0 = -t$, $x_i = x^i$, $i = 1,\ldots,n$. 我们采用记号 $\partial_\mu = \frac{\partial}{\partial x^\mu}$ 与 $\partial_t = \frac{\partial}{\partial x^\mu}$(原文如此). 记 $\Sigma_{t_0}$ 为类空超平面 $t = t_0$. 波动算子定义为 $\Box = m^{\alpha\beta}\partial_{\alpha\beta} = -\partial_t^2 + \sum_i\partial_i^2$.
下面我们回顾方程 (1) 在 $t=0$ 处满足初值问题的解的基本已知估计:
\begin{equation}
\phi(0,x)=f(x),\qquad \partial_{t}\phi(0,x)=g(x).\end{equation}

为方便起见, 记 $\phi[0] = (f, D^{-1}g)$, 其中 $D^{-1}$ 为象征 $|\xi|^{-1}$ 的拟微分算子.
\begin{proposition}[能量恒等式、不等式]\label{prop:1.1}
方程 (1)、(2) 的解满足
\begin{equation}
E[\partial\phi](t)=E[\partial\phi](0),\end{equation}
其中
\begin{equation}
E[\partial\phi](t)=\int_{\Sigma_{t}}\bigl(|\partial_{t}\phi|^{2}+\sum_{i}|\partial_{i}\phi|^{2}\bigr)\,d x.\end{equation}
由此可得能量不等式: 对任意 $s \geq 0$,
\begin{equation}
\|\partial\phi(t)\|_{H^{s}(\mathbb{R}^{n})}\leq\|\partial\phi(0)\|_{H^{s}(\mathbb{R}^{n})}.\end{equation}
\end{proposition}

\begin{proof}
能量恒等式既可由只涉及分部积分的几何方法证明, 也可利用 Plancherel 公式与傅里叶表示公式, 通过傅里叶方法证明:
\begin{equation}
\phi(t,x)=(2\pi)^{-n}\int e^{i x\cdot\xi}\Bigl(\cos t|\xi|\,\hat f(\xi)+\frac{\sin t|\xi|}{|\xi|}\,\hat g(\xi)\Bigr)\,d\xi.\end{equation}
更高阶的能量不等式基于 $\Box$ 与 $\partial_{\mu}$ 的可交换性.\end{proof}

\begin{remark}\label{rem:1.2}
标准的 Sobolev 嵌入 $H^s(\mathbb{R}^n) \subset L^\infty(\mathbb{R}^n)$(当 $s > \frac{n}{2}$)使我们无需利用显式表示即可得到方程 (1) 解的 $L^\infty$ 估计. 这一手续在与局部存在定理的联系中可以推广到非线性方程.\end{remark}

\begin{proposition}[色散不等式]\label{prop:1.3}
方程 (1)、(2) 的解满足
\begin{equation}
\|\phi(t)\|_{L^{\infty}}\leq c\,t^{-\frac{n-1}{2}}\|D^{\frac{n+1}{2}}\phi[0]\|_{L^{1}}.\end{equation}
\end{proposition}

\begin{remark}\label{rem:1.4}
事实上 (6) 并不十分准确; 如果把左边的 $L^\infty$ 范数换成 $\mathrm{BMO}$ 范数, 或者把右边的 $L^1$ 范数换成 Hardy 范数 $\mathcal{H}^1$, 则正确的估计成立. 不过, 若数据 $\phi(0) = f$、$\partial_t \phi(0) = g$ 的傅里叶变换支在某个二进球壳 $\frac{\lambda}{2} \leq |\xi| \leq 2\lambda$(其中某个固定的 $\lambda \in 2^{\mathbb{N}}$)中, 则 (6) 如所书写是成立的.\end{remark}

\begin{proof}
直到最近, (6) 唯一稳健的证明都基于对表示式 (5) 的驻相法. 在奇数维, 利用解的角均值表示也可以证明色散估计的一种相关形式. 我们后面将讨论一种避免任何表示公式的 (6) 的推导.\end{proof}

\begin{remark}\label{rem:1.5}
色散不等式提供了两类信息. 第一类涉及当 $t \to \infty$ 时 $\|\phi(t)\|_{L^\infty}$ 的精确衰减率; 第二类涉及对 $t > 0$ 时 $\|\phi(t)\|_{L^\infty}$ 的正则性性质. 就改进的正则性而言, 与 Sobolev 嵌入 $L^\infty(\mathbb{R}^n) \subset W^{1,n}(\mathbb{R}^n)$ 相比, (6) 在 $t > 0$ 时多获得了 $\frac{n-1}{2}$ 阶导数.\end{remark}

众所周知, 就渐近行为而言, (6) 在应用于非线性波动方程时并不十分有用. 推导波动方程解渐近性质的一个更有效的方法, 基于由交换向量场法得到的广义能量估计, 以及全局 Sobolev 不等式. 下面我们极快地介绍这一观点, 并在第一讲的补充中更系统地展开. 我们回顾用交换向量场法推导上述衰减率估计. 其思想是将能量恒等式 (3) 与交换向量场以及经典 Sobolev 不等式的全局形式结合使用.
Minkowski 时空 $\mathbb{R}^{n+1}$ 装备有一族 Killing 与保形 Killing 向量场: 平移 $T_{\mu} = \partial_{\mu}$, Lorentz 旋转 $O_{\mu\nu} = x_{\mu} \partial_{\nu} - x_{\nu} \partial_{\mu}$, 伸缩 $S = t \partial_{t} + x^{i} \partial_{i}$, 以及反平移 $K_{\mu} = -2x_{\mu}S + \langle x, x\rangle \partial_{\mu}$. 这里 $x^{\mu}$ 表示标准变量 $x^{0} = t, x^{1}, \ldots, x^{n}$, 而 $x_{\mu} = m_{\mu\nu} x^{\mu}$. Killing 向量场 $T_{\mu}$ 与 $O_{\mu\nu}$ 与 $\Box$ 交换, 而 $S$ 以保持解空间的意义与 $\Box$ 交换: 由 $[\Box, S] = 2\Box$ 知 $\Box \phi = 0$ 蕴含 $\Box S \phi = 0$. 我们把算子 $O_{\mu\nu}$ 分解为角旋转算子 $(ij)$: $O = x_{i} \partial_{j} - x_{j} \partial_{i}$ 与推进 $(i)$: $L = x_{i} \partial_{t} + t \partial_{i}$, 其中 $i, j, k = 1, \ldots, n$. 回顾 (4) 中的能量表达式; 基于上述交换性质我们定义如下"广义能量"
\begin{equation}
E_{k}[\partial\phi]=\sum_{X_{i_{1}},\ldots,X_{i_{j}}}E^{2}\bigl[\partial X_{i_{1}}X_{i_{2}}\ldots X_{i_{j}}\phi\bigr],\end{equation}
其中求和取遍 $0 \leq j \leq k$ 以及所有 Killing 向量场 $T, L_{\mu\nu}$ 与伸缩向量场 $S$. 交换向量场法的要点在于, 量 $E_k\;(k \geq 1)$ 被方程 (1) 的解守恒. 因此, 若对所有 $0 \leq k \leq s$ 数据 $f, g$ 满足
\begin{equation}
\int(1+|x|)^{2k}\Bigl(|\nabla^{k+1}f(x)|^{2}+|\nabla^{k} g(x)|^{2}\Bigr)\,d x\leq C_{s}<\infty,\end{equation}
则对所有 $t$ 有 $E_s[\partial\phi](t) \leq C_s$. 方程 (1) 的所需衰减估计现在可由下面的 Sobolev 不等式的全局形式(见 \cite{klai27})推出:
\begin{proposition}[衰减估计, 见 {\cite{klai27}}]\label{prop:1.6}
设 $\phi$ 为 $\mathbb{R}^{n+1}$ 中任一函数, 使得对某整数 $s > \frac{n}{2} + 1$ 有 $E_s[\phi]$ 有限. 则对所有 $t > 0$,
\begin{equation}
|\partial\phi(t,x)|\lesssim(1+t+|x|)^{-\frac{n-1}{2}}(1+|t-|x||)^{-\frac{1}{2}}E_{s}[u].\end{equation}
因此, 若 (1) 中的数据 $f, g$ 对 $0 \leq k \leq s$ 满足 (8) 且某 $s > \frac{n}{2}$, 则对所有 $t \geq 0$,
\begin{equation}
\|\partial\phi(t,\cdot)\|_{L^{\infty}}\lesssim\frac{1}{(1+t+|x|)^{\frac{n-1}{2}}(1+|t-|x||)^{\frac{1}{2}}}.\end{equation}
\end{proposition}

\begin{remark}\label{rem:1.7}
显然, 这个证明纯为几何$^{1}$ 的估计蕴含了色散不等式 (6) 给出的衰减性质. 事实上, 它在波区 $|x| \sim t$ 之外提供了更多信息, 与线性方程 $\Box \phi = 0$ 预期的传播性质十分吻合. 另一方面, 由于 (9) 实质上是 Sobolev 不等式的全局形式, 似乎\propref{prop:1.6} 的估计与 (6) 改进的正则性特征无关. 然而这并不正确, 我们将在下一个命题中看到.\end{remark}

\begin{remark}\label{rem:1.8}
交换向量场法(\cite{klai27}, 亦见 \cite{klai46})最初是作为 \cite{klai17} 几乎整体存在结果的另一种大为简化的证明而发展的. 该方法非常灵活, 适用于证明维数 $n>3$ 的整体存在性以及维数 $n=3$ 对一大类非线性波动方程(包括完全非线性的)在小初值下的几乎整体存在性. 它还与零条件(见\propref{prop:1.22} 后的讨论)结合, 证明低维 $n=3,2$ 的整体存在性. 向量场法还被许多其他作者发展, 特别是 H. Lindblad \cite{klai48} 与 T. Sideris \cite{klai63,klai64,klai65}, 亦见 Hörmander \cite{klai16} 及其中的参考文献. 向量场法的一个重要修正出现在 \cite{klai12} 中, 其中向量场必须被修改. 这将在第三讲中一定程度上讨论.\end{remark}

\begin{proposition}[见 {\cite{klai30}}]\label{prop:1.9}
交换向量场法蕴含色散不等式 (6).\end{proposition}

\begin{proof}
不失一般性可设 $\partial_t \phi = g = 0$, 且 $f = \phi(0)$ 的傅里叶变换支在某个壳 $\frac{\lambda}{2} \leq |\xi| \leq 2\lambda$($\lambda \in 2^{\mathbb{N}}$)中. 通过简单的伸缩论证, 实际上可设 $\lambda = 1$. 由于 $\phi$ 关于空间变量 $x$ 的傅里叶变换 $\hat{\phi}$ 也支在同一壳中, 只需对 $\nabla \phi$ 或 $\nabla^k \phi$ 证明估计即可. 接下来用并排半径为有限、互不重叠的盘 $D_I$ 覆盖 $\mathbb{R}^n$, 这些盘以整数坐标点 $I \in \mathbf{Z}^n$ 为中心, 使得每个 $D_I$ 至多与有限个 $c_n$ 个盘 $D_J$ 相交, 其中 $c_n$ 只依赖于维数 $n$. 考虑光滑单位分解 $(\chi_I)_{I\in\mathbf{Z}^n}$, 满足 $\operatorname{supp}\chi_I\subset D_I$ 且各 $\chi_I$ 为正. 显然可以安排使得对所有 $k$,
\[
\sum_{I\in\mathbf{Z}^{n}}|\nabla^{k}\chi_{I}(x)|\leq C_{k,n}
\]
对 $x \in \mathbb{R}^n$ 一致成立. 现在令 $f_I = \chi_I \cdot f$, $\phi_I$ 为对应的方程 (1) 的解, 数据 $\phi_I(0) = f_I, \partial_t \phi_I(0) = 0$. 显然 $f = \sum_I f_I$, $\phi = \sum_I \phi_I$. 只需证明对所有 $I$,
\[
\|\nabla^{k}\phi_{I}(t)\|_{L^{\infty}}\leq C_{n,k}(1+t)^{-\frac{n-1}{2}}\sum_{j=0}^{n+k+1}\|D^{j}f_{I}\|_{L^{1}},
\]
其中常数 $C_{n,k}$ 只依赖于 $n$ 与 $k$. 确实, 若 (12) 成立, 则易得 $\|\nabla^{k}\phi(t)\|_{L^\infty} \leq C_{n,k}(1+t)^{-\frac{n-1}{2}} \sum_{j=0}^{n+k+1} \|\sum_{I} \nabla^{j}\chi_{I}\|_{L^\infty} \|f\|_{L^{1}}$, 从而由 (11) 得 $\|\nabla^{k}\phi(t)\|_{L^\infty} \leq C_{n,k}(1+t)^{-\frac{n-1}{2}} \|f\|_{L^{1}}$.
于是只需验证 (12). 经空间平移不妨设 $I = 0$. 把\propref{prop:1.6} 用于 $\psi = \nabla \phi_0$, 对 $s_*$ 为严格大于 $\frac{n}{2} + 1$ 的首个整数, 得
\[
\begin{aligned}
\|\psi(t)\|_{L^{\infty}}&\leq c(1+t)^{-\frac{n-1}{2}}E_{s_{*}}[\phi_{0}](t)\\
&\leq c(1+t)^{-\frac{n-1}{2}}E_{s_{*}}[\phi_{0}](0).\end{aligned}
\]
由于 $\phi_{0}$ 的支集包含于以原点为中心、半径为 1 的球内, 有
\[
E_{s_{*}}[\phi_{0}](0)\leq C_{n}\sum_{j=0}^{s_{*}+1}\|D^{j}f_{0}\|_{L^{2}}.
\]
最后, 根据 $\mathbb{R}^n$ 中的标准 Sobolev 不等式 $\|f\|_{L^2} \leq c\|\nabla^{\frac{n}{2}}f\|_{L^1}$, 得到
\[
\|\psi(t)\|_{L^{\infty}}\leq c(1+t)^{-\frac{n-1}{2}}\sum_{j=0}^{n+2+1}\|D^{j}f_{0}\|_{L^{1}},
\]
如所求.\end{proof}

在能量不等式与色散不等式之间作插值, 得到所谓的 Strichartz--Brenner 结果:
\begin{equation}
\|\phi(t)\|_{L^{r}}\leq c\,|t|^{-\gamma(r)}\|\nabla^{\sigma}\partial\phi(0)\|_{L^{r^{\prime}}},\end{equation}
其中 $\gamma(r) = (n-1)\bigl(\frac{1}{2}-\frac{2}{r}\bigr)$, $\frac{1}{r}+\frac{1}{r'}=1$, $r\geq2$, 且伸缩条件为 $\frac{n}{r}=-\gamma(r)-\sigma-1+\frac{n}{r'}$. 通过标准的 $TT^{*}$ 论证、Hardy--Littlewood--Sobolev 不等式以及 Littlewood--Paley 理论的应用, 得到广义 Strichartz 型不等式.
\begin{definition}\label{def:1.10}
一对正实数 $(q, r)$ 称为波容许(wave admissible), 如果
\[
\frac{2}{q}\leq\gamma(r)=(n-1)\Bigl(\frac{1}{2}-\frac{2}{r}\Bigr),\qquad q\geq2,\qquad(q,r,n)\neq(\infty,1,3).
\]
\end{definition}

\begin{proposition}[Strichartz 型不等式]\label{prop:1.11}
初值问题 (1)--(2) 的解满足
\[
\|\phi\|_{L_{t}^{q}L_{x}^{r}([0,T]\times\mathbb{R}^{n})}
\leq c\,\|\phi[0]\|_{H^{\sigma}},\qquad \sigma=n\Bigl(\frac{1}{2}-\frac{2}{r}\Bigr)-\frac{1}{q}.
\]
\end{proposition}

\begin{remark}\label{rem:1.12}
Strichartz 本人只证明了各向同性情形 $q = r$ 的估计. 其余估计归于 Pecher、Ginibre--Velo、Lindblad--Sogge 与 Keel--Tao. 关于这些估计的最终版本与精确参考文献, 我们参阅 Keel--Tao \cite{klai20}.\end{remark}

\begin{remark}\label{rem:1.13}
Strichartz 型不等式在非线性波动方程理论的许多最新进展中起着关键作用. 注意, 在固定频率下由能量恒等式与色散不等式推出 (13) 的步骤相当柔和, 它们可追溯到 Duhamel 原理与初值问题的唯一性$^{2}$. 两者都适用于一般变系数线性波动方程, 且只要求系数有很少的正则性. 因此 Strichartz 型不等式的主要构件是 (3) 与 (6).\end{remark}

迄今为止讨论的所有估计都只涉及波动方程 (1) 的单个解. 它们在应用于非线性项含导数的波动方程时价值有限. 过去十年中发展出一种新观点: 直接估计涉及 (1) 的多个解的二次甚至多次线性表达式, 可获得远为灵活的估计.
\begin{proposition}[双线性 Strichartz, 见 {\cite{klai47}}]\label{prop:1.14}
设 $\phi, \psi$ 为 (1) 的两个解. 设 $n \geq 2$ 且 $(q, r)$ 为波容许. 则
\[
\|D^{\sigma}(\phi\psi)\|_{L_{t}^{q/2}L_{x}^{r/2}}\lesssim\|\phi[0]\|_{\dot{H}^{s_{1}}}\|\psi[0]\|_{\dot{H}^{s_{2}}},
\]
其中 $\sigma, s_{1}, s_{2}$ 满足
\[
0<\sigma<n-\frac{2n}{r}-\frac{4}{q},\qquad
0<s_{1},s_{2}<\frac{n}{2}-\frac{n}{r}-\frac{1}{q},\qquad
s_{1}+s_{2}+\sigma=n-\frac{2n}{r}-\frac{2}{q}.
\]
\end{proposition}

\begin{proof}
证明基于如下色散不等式的微局部形式:\end{proof}

\begin{proposition}[改进的色散不等式, 见 {\cite{klai47}}]\label{prop:1.15}
考虑 (1)(2) 的一个解, 其初数据的傅里叶变换支在半径为小 $\mu$ 的球内, 该球包含于单位二进环 $\frac{1}{2} \leq |\xi| \leq 2$ 中. 则
\[
|\phi(t)|_{L^{\infty}}\leq c\,\mu\, t^{-\frac{n-1}{2}}\|\phi[0]\|_{L^{1}}.
\]
\end{proposition}

改进的色散不等式易由标准驻相法得到. 我们至今还没有满意的几何推导.
现在考虑更一般的双线性估计, 形如
\[
\left\|D^{\gamma}D_{+}^{\gamma+}D_{-}^{\gamma-}\left(\phi\psi\right)\right\|_{L_{t}^{q/2}L_{x}^{r/2}}\lesssim\|\phi[0]\|_{\dot{H}^{s_{1}}}\|\psi[0]\|_{\dot{H}^{s_{2}}},
\]
其中 $\phi$ 与 $\psi$ 满足 (1). 这里 $D^\gamma$、$D_{+}^{\gamma+}$、$D_{-}^{\gamma-}$ 为象征 $|\xi|^\gamma$、$|\tau| + |\xi||^{\gamma+}$ 与 $|\tau| - |\xi||^{\gamma-}$ 的乘子算子. 在 $q, r = 4$ 的情形, 所有此类估计都已知道. 下述定理的特例先后出现在 \cite{klai31}、以及 \cite{klai37,klai44,klai39} 中. 完整解最近由 Foschi--Klainerman \cite{klai14}(亦见 \cite{klai69})给出.
\begin{proposition}[双线性估计, 见 {\cite{klai14}}]\label{prop:1.16}
设 $n \geq 2$ 且 $\gamma, \gamma_{-}, \gamma_{+}, s_1, s_2 \in \mathbb{R}$. 估计
\[
\left\|D^{\gamma}D_{+}^{\gamma+}D_{-}^{\gamma-}\left(\phi\psi\right)\right\|_{L_{t}^{2}L_{x}^{2}}\lesssim\left\|\phi[0]\right\|_{\dot{H}^{s_{1}}}\left\|\psi[0]\right\|_{\dot{H}^{s_{2}}}
\]
对 (1) 的所有初数据成立, 当且仅当下列条件满足:
\[
\gamma+\gamma_{+}+\gamma_{-}=s_{1}+s_{2}-\frac{n-1}{2},\qquad
\gamma_{-}\geq-\frac{n-3}{4},\qquad
\gamma>-\frac{n-1}{2},
\]
\[
s_{i}\leq\gamma_{-}+\frac{n-1}{2}\quad (i=1,2),\qquad
s_{1}+s_{2}\geq\frac{1}{2},
\]
\[
\left(s_{i},\gamma_{-}\right)\neq\left(\frac{n+1}{4},-\frac{n-3}{4}\right)\quad (i=1,2),\qquad
\left(s_{1}+s_{2},\gamma_{-}\right)\neq\left(\frac{1}{2},-\frac{n-3}{4}\right).
\]
\end{proposition}

这些估计与零二次型密切相关.
\begin{definition}\label{def:1.17}
设 $\phi, \psi$ 为 $\mathbb{R}^{n+1}$ 上的两个光滑函数. 定义
\[
Q_{0}(\phi,\psi)=m^{\alpha\beta}\partial_{\alpha}\phi\partial_{\beta}\psi,\qquad
Q_{i j}(\phi,\psi)=\partial_{i}\phi\partial_{j}\psi-\partial_{i}\psi\partial_{j}\phi,\qquad
Q_{0i}(\phi,\psi)=\partial_{i}\phi\partial_{t}\psi-\partial_{i}\psi\partial_{t}\phi.
\]
\end{definition}

\begin{definition}\label{def:1.18}
若 $\phi, \psi$ 为 $\mathcal{S}'(\mathbb{R}^{n+1})$ 中的分布, 其时空傅里叶变换 $\widehat{\phi}, \widehat{\psi}$ 为函数. 记 $\phi \preceq \psi$ 当且仅当 $|\widehat{\phi}| \leq C\widehat{\psi}$.\end{definition}

\begin{lemma}[见 {\cite{klai44}}]\label{lem:1.19}
下列成立:
\[
Q_{0}(\phi,\psi)\precsim D_{+}D_{-}(\phi\psi),\qquad
Q_{i j}(\phi,\psi)\prec{\cal D}^{\frac{1}{2}}{\cal D}_{-}^{\frac{1}{2}}({\cal D}^{\frac{1}{2}}\phi\,{\cal D}^{\frac{1}{2}}\psi).
\]
\end{lemma}

\begin{remark}\label{rem:1.20}
对 $Q_{0i}$ 也可推出类似公式. 可结合\propref{prop:1.16} 的双线性估计与上面的引理, 导出零二次型的非平凡 $L^{2}$ 时空估计.\end{remark}

\begin{remark}\label{rem:1.21}
尽管\propref{prop:1.16} 的双线性估计中正则性的获得是近十年的进展, 但零二次型对非线性方程应用的重要性此前已在衰减获得(gain of decay)的联系中为人所知. 下面是一个典型事实.\end{remark}

\begin{proposition}[零形式的衰减, 见 {\cite{klai28}}]\label{prop:1.22}
设 $\phi,\psi$ 为 (1) 的两个解, 并设其初数据对充分大的 $s$ 满足加权 Sobolev 范数估计 (8). 则对任意零形式 $Q_{0},Q_{ij},Q_{0i}$, 有衰减估计
\[
\|Q(\phi,\psi)(t)\|_{L^{2}(\mathbb{R}^{n})}\leq C\, t^{-n}
\]
(当 $t \to \infty$).\end{proposition}

证明是\propref{prop:1.6} 衰减估计以下改进的直接推论.
\begin{proposition}[剥离衰减性质]\label{prop:1.23}
设 $\phi$ 为 $R^{n+1}$ 中使得对某整数 $s > \frac{n}{2} + 2$ 有 $E_s[\phi]$ 有限的函数. 令 $E_\pm = \partial_r \pm \partial_r$(其中 $\partial_r = \sum_i \frac{x^i}{|x|}\partial_i$), 且 $E_i = 0$, $E_A, E_B \geq \delta_{AB}$. 记 $|\nabla \phi|^2 = \sum_{i=1}^{n-1} E_A(\phi)^2$. 我们有
\[
\|E_{+}\phi(t)\|_{L^{\infty}}\lesssim\left(\frac{1}{t}\right)^{\frac{n+1}{2}},
\]
\[
\|\bar{\nabla}\phi(t)\|_{L^{\infty}}\lesssim\left(\frac{1}{t}\right)^{\frac{n+1}{2}},
\]
\[
\|\left(1+|u|\right)E_{-}\phi(t)\|_{L^{\infty}}\lesssim\left(\frac{1}{t}\right)^{\frac{n-1}{2}}.
\]
\end{proposition}

\begin{remark}\label{rem:1.24}
\propref{prop:1.22} 中零二次型改进衰减的应用, 见 \cite{klai28}, 用于证明 $\mathbb{R}^{3+1}$ 中一大类满足零条件的拟线性波动方程的小初值整体存在性. 证明 \cite{klai28} 中结果所用的向量场技巧被 Sideris 在 \cite{klai63}、\cite{klai64}、\cite{klai65} 中推广到 Wave Maps 与非线性弹性力学.\end{remark}

至此我们只考虑了 $q = r = 4$ 型 (16) 的一般双线性估计. 对于"波容许"指数 $q, r$ 与 $\gamma_{-}$, 当 $s_{1} = s_{2}$ 时\propref{prop:1.14} 给出的估计是最优的. 存在一个一般双线性猜想, 断言在指数 $\gamma, \gamma_{+}, \gamma_{-}, s_1, s_2$ 的适当限制下估计 (16) 成立. T. Tao \cite{klai68} 利用 T. Wolff \cite{klai76} 在其锥限制猜想近期重要工作中发展的思想, 证明了该猜想的一些非常重要情形. 不过, 我想以这一近期进展中一个重要而非常简单的要素来结束第一讲. 这就是一般波动可分解为管(tubes)的思想.
\begin{definition}\label{def:1.25}
称波 $\phi$ 为 $\pm$ 波, 如果其初数据满足
\[
\partial_{t}\phi(0)=\pm i\sqrt{-\Delta}\,\phi(0),
\]
或在傅里叶表示中 $\partial_t\hat\phi(\theta,\xi) = \pm i|\xi|\hat\phi(\theta,\xi)$. 注意任意波可按公式
\[
\phi=\phi^{+}+\phi^{-}
\]
分解为 $\pm$ 波, 其中 $\phi^{\pm}$ 的初数据为 $\phi^{\pm}[0] = P^{\pm}\phi[0]$,
\[
P_{\pm}\phi[0]:=\Bigl(\frac{\phi(0)\mp i(-\Delta)^{-1/2}\phi_{t}(0)}{2},\frac{\phi_{t}(0)\pm i(-\Delta)^{1/2}\phi(0)}{2}\Bigr).
\]
\end{definition}

\begin{remark}\label{rem:1.26}
注意解 $\phi^+, \phi^-$ 由其数据唯一确定. 换言之, 为定义它们只需分解 $\phi$ 的数据. 不过它们可用傅里叶变换表示为
\[
\phi^{\pm}(t,x)=\int e^{\pm i t|\xi|}e^{i x\cdot\xi}\widehat{P^{\pm}\phi[0]}(\xi)\,d\xi.
\]
\end{remark}

\begin{proposition}\label{prop:1.27}
考虑 $\phi = \phi^{-}$ 为一支减波, 其 $\phi^{-}[0]$(本质上)支在球 $B = B(0, \Lambda_0)$ 中, $\Lambda_0 > 1$, 且其傅里叶变换支在区域 $\{\xi: \frac{1}{2} \leq \xi_1 \leq 2,\ |\xi'| \leq \frac{1}{\Lambda_0}\}$ 中, 其中 $|\xi'| = \sqrt{\xi_2^2 + \cdots + \xi_n^2}$. 对任意 $t \geq 0$ 考虑区域 $T^+(t, \Lambda) = \{x: |x| \leq t + \Lambda_0;\ |x_1 - t|^2 + |x'|^2 \leq \Lambda^2\}$. 则对所有满足 $\Lambda_0 \leq \Lambda \leq t \leq \Lambda_0^2$ 的 $t$:
\[
\|\partial\phi^{-}(t)\|_{L^{2}(\Sigma_{t}\backslash T^+(t,\Lambda))}\leq C_{N}\Bigl(\frac{\Lambda_{0}}{\Lambda}\Bigr)^{N}E(\phi)^{\frac{1}{2}}.
\]
\end{proposition}

\begin{proof}
该命题的标准证明基于积分表示公式 (20) 中的简单分部积分. 也可给出基于交换向量场 \cite{klai43} 的纯几何证明.\end{proof}

上述命题的向量场法证明, 结合\propref{prop:1.22} 的零估计以及由 T. Wolff \cite{klai76} 与 T. Tao \cite{klai68} 首创的"尺度归纳"(induction on scales)重要思想, 使我们$^{4}$ 能给出如下双线性估计的纯几何证明:
\begin{proposition}\label{prop:1.28}
设 $\phi, \psi$ 为 (1) 的两个解, $Q$ 为零形式 $Q_{0}, Q_{ij}$ 或 $Q_{0i}$ 之一. 则
\[
\|Q(\phi,\psi)\|_{L^{2}(\mathbb{R}^{n+1})}\lesssim\|D\phi[0]\|_{L^{2}(\mathbb{R}^{n})}\|D^{\frac{n+1}{2}}\phi[0]\|_{L^{2}(\mathbb{R}^{n})}.
\]
\end{proposition}

这一估计是\propref{prop:1.16} 的特例, 最早由 Machedon 与我和 \cite{klai31} 借助 $Q(\phi, \psi)$ 的时空傅里叶变换大量依赖地证明. 它已被用于证明 Yang--Mills \cite{klai34}、Maxwell--Klein--Gordon \cite{klai33} 与临界幂 Yang--Mills--Higgs \cite{klai19} 的有限能量整体存在性.
